
A new polyconvex electro-thermo-visco-hyperelastic constitutive framework has been presented for EAPs operating under finite deformation. The proposed thermal constitutive structure overcomes the intrinsic limitations associated with constant or purely temperature-dependent specific heat formulations. In particular, satisfaction of the Third Law is achieved without compromising polyconvexity, guaranteeing a strictly positive heat capacity for arbitrary admissible deformation states and preserving the stability characteristics of the constitutive model through positive wave speeds.

The equilibrium/non-equilibrium decomposition provides a consistent framework for incorporating finite-strain viscoelasticity while maintaining a transparent separation between reversible and dissipative mechanisms. Beyond the isotropic formulation presented here, this constitutive architecture naturally accommodates anisotropic extensions through appropriate structural tensors introduced independently within the equilibrium and non-equilibrium free-energy contributions.

Calibration against available thermo-mechanical, electro-mechanical and coupled electro-thermo-mechanical experiments demonstrates that the proposed model accurately reproduces a broad range of observed responses despite the significant variability present in published datasets. Finite element simulations further illustrate the role of viscous dissipation as an internal heat source and quantify its influence on the coupled electro-thermo-mechanical response under large deformation.

The present framework establishes a constitutive basis in which thermodynamic admissibility, mathematical stability and computational robustness are addressed simultaneously. As such, it provides a general platform for future developments including anisotropic electro-active polymers, electro-chemo-mechanical coupling, damage, fatigue and other multiphysics extensions within a unified setting.



% This section summarizes the primary findings and contributions of the present work. To provide a structured overview, the conclusions are categorized into three distinct areas: the theoretical development of the thermo-electro-mechanical constitutive model, the rigorous parameter identification and calibration strategy, and the global implications of the proposed computational framework.


% \subsection{Thermodynamically consistent constitutive model}




% \subsection{Calubration of the constitutive model}

% First of all, in order to perform a consistent calibration of the constitutive model presented in \ref{sec:constitutive equations}, only a few experimental data has been selected among the literature (see tables \ref{tab:experiment_sources} and \ref{tab:calibration_steps}) because of the big discrepancies of reported values for the same experiment under the same conditions among different laboratories, and the need to isolate the uncertainty from data acquisition and the uncertainty inherent to the limitations of the constitutive model. Even thought, data from three different research groups has been combined for the calibration and the overall result has proven to be accurate. Key findings from the calibration and validation procedures include:

% \begin{itemize}
%   \item The calibration of the thermally-coupled volumetric energy, even if the variation of the specific heat capacity is almost linear \cite{Dippel2015}, has demonstrated to be consistent in terms of energy, entropy and specific heat capacity. Furthermore, an indirect characterization of the bulk modulus $\kappa_R$ has been performed and it is in good agreement with standard values.
%   \item Regarding the deviatoric response of the material, it has proven to be more accurate to perform a staggered characterization of the long term component and the viscous branches. From the numerical optimizations, different optimal parameters had been observed if a monolithic calibration would be performed. Additionally, models that depend solely on the first invariant $I_1$, do not seem to accurately capture the the response of the material, for that reason, a nonlinear Mooney-Rivlin model has been chosen. For the sake of simplicity, the viscous branches are modelled with simple neo-Hookean models, however, the choice of a more complex model would easily enhance the fitting, at the cost of increasing the number of parameters.
%   \item The observation of the error metrics such as the sensitivity index for each calibrated parameter had allowed a quantitative criterion the choose the number of viscous branches. Three branches provide the best fitting, even two viscous branches with a neo-Hookean model provide good results, adding a fourth branch would result into an overfitted model.
%   \item With regard to the thermoelasticity coupling, each component of the free energy density has been modelled with an independent thermal law, a nonlinear softening has been proposed, as a direct extension of the model presented in \cite{BonetGil2026}. Due to statistical observability, all the viscous branches have been modelled with the same thermal function.
%   \item Finally, two types of validation have been performed. The calibrated model was used to reproduce the thermo-electro-mechanical experiments presented in an independent publication, offering a blind validation with a fitting $R^2=93\%$. additionally, the specific heat capacity has been computed for different combinations of stretch and temperature, to validate the positivity of the specific heat and fulfilment of the third law of the thermodynamics.
% \end{itemize}

% \subsection{General conclusions}

